\documentclass{article}
\usepackage{graphicx} % Required for inserting images

\title{Sustainability-Aware OS Scheduling for Servers}
\author{Ohad Eitan}
\date{June 2026}

\begin{document}

\maketitle

\section{Introduction}

We propose extending operating-system schedulers on servers to optimize for
sustainability, not just throughput or fairness. The core observation is that the
cost of running a job is not constant: electricity prices and the carbon intensity
of the grid vary throughout the day, and a server's available capacity changes as
its utilization fluctuates. A scheduler that is aware of these factors can shift
work toward cheaper, cleaner, less-congested time windows.

\section{Core Idea}

We divide server time into discrete blocks (which may be fixed or dynamically
sized). Each block has an associated cost derived from:

\begin{itemize}
    \item \textbf{Electricity price} at that time of day, which is non-uniform
    across the day.
    \item \textbf{Carbon intensity} of the grid at that time (to minimize
    CO\textsubscript{2}).
    \item \textbf{Current utilization} of the server --- the remaining capacity
    below 100\%.
\end{itemize}

The result behaves like a market: when load is low, running a job is cheaper; when
the server is busy or electricity is expensive, it costs more. Users are thereby
incentivized to schedule flexible work into off-peak, low-carbon windows.

\section{Scheduling Guidelines}

\begin{enumerate}
    \item No job is killed once it has been submitted and admitted.
    \item Minimize CO\textsubscript{2} emissions across all scheduled work.
    \item Price tracks load and time of day --- lower load means lower cost, like
    a stock market.
    \item Encourage longer, batched jobs. We want users to consolidate work so it
    can run as a group, improving efficiency.
\end{enumerate}

\section{User Submission}

A user submits a job request with the following parameters:

\begin{itemize}
    \item \textbf{Preferred time of day} --- when the user would like the job to
    run.
    \item \textbf{Estimated length} --- the job's expected duration, from the
    user's perspective.
    \item \textbf{Accuracy score} --- a measure of how accurate this user's
    previous length estimates have been (their track record).
\end{itemize}

The accuracy score feeds into a penalty/reward (``punishment'') mechanism: users
with a history of accurate estimates are treated more favorably, while consistently
inaccurate estimates incur a penalty. This discourages users from gaming the system
with misleading job lengths.

\section{Scheduler Inputs}

The scheduler computes its decisions from:

\begin{itemize}
    \item Current load and utilization of the server.
    \item Cost of a block and length of a block.
    \item Time of day.
    \item Cost per user per task (with the goal of nudging users toward longer,
    grouped jobs).
    \item User's historical estimate accuracy.
\end{itemize}

\end{document}
